\chapter{Felhasználói dokumentáció}
\label{ch:user}

\section{A program célja} 
Ez a fejezet a rendszer használatát mutatja be a felhasználó szemszögéből. Ismerteti az alkalmazás elérésének és telepítésének módját, bemutatja a főbb funkciókat, valamint segítséget nyújt azok használatához. A fejezet végén az ismert korlátozások, az elérhetőségi információk és a hibaelhárítással kapcsolatos tudnivalók találhatók.
\mbox{}

\section{Rendszerkövetelmények és indítás}
\subsection{Rendszerkövetelmények}
A rendszer online eléréséhez nincs különösebb hardverigény.
Támogatott böngészők: Google Chrome, Microsoft Edge és Mozilla Firefox.
Az online változat használatához internetkapcsolat és egy támogatott böngésző szükséges. A helyi futtatáshoz ezen felül a 2.2.3 fejezetben ismertetett szoftverek telepítése is szükséges.

\mbox{}

\subsection{Az alkalmazás elérése}
Az alkalmazás GitHub Pages-en van közzétéve, és az alábbi linken érhető el:
\begin{center}
  \url{https://szaboanna21.github.io/thesis/#/}
\end{center}

\mbox{}

\subsection{Helyi futtatás}
A felület internetkapcsolat nélkül, lokálisan is futtatható a felhasználó számítógépén.
A futtatáshoz az alábbi szoftverek szükségesek:
\begin{itemize}
  \item \textbf{Node.js 24} (\url{https://nodejs.org}), amely lehetővé teszi a JavaScript alapú alkalmazások futtatását a számítógépen. 
  \item \textbf{pnpm 12.6}: \texttt{npm install -g pnpm@12.6.0}, amely automatikusan telepíti és kezeli az alkalmazás működéséhez szükséges csomagokat.

Hogyha, ezek már telepítve vannak, az alkalmazás az alábbi lépésekkel indítható el:

\begin{enumerate}
  \item A projekt forráskódjának letöltése: 
  \begin{lstlisting} 
    git clone https://github.com/SZaboAnna21/thesis.git
  \end{lstlisting} 

  \item Belépés a projekt könyvtárába: 
  \begin{lstlisting} 
    cd thesis/Fotoc-helper
  \end{lstlisting} 
  \item Az alkalmazás működéséhez szükséges csomagok telepítése:
  \begin{lstlisting} 
    pnpm install
  \end{lstlisting} 
  Ez a lépés automatikusan letölti a program futtatásához szükséges kiegészítő könyvtárakat.
  \item Indítás: 
  \begin{lstlisting} 
    ppnpm --filter web dev
  \end{lstlisting}  
  \item A sikeres indítást követően az alkalmazás a böngészőből a \texttt{http://localhost:5173} címen érhető el.
\end{enumerate}

\mbox{}
A helyi futtatás elsősorban fejlesztői vagy tesztelési célokat szolgál. 
Azoknak a felhasználóknak, akik kizárólag használni szeretnék az alkalmazást, 
elegendő az online változat megnyitása a 2.2.2 fejezetben szereplő címen.

\section{A felület áttekintése}

\subsection{Főoldal és navigáció}
A fenti „Az alkalmazás elérése” szekcióban bemutatott hivatkozás a rendszer főoldalára navigálja a felhasználót.
 A főoldal rövid ismertetőt tartalmaz az alkalmazás céljáról és használatáról, valamint innen indulva érhető el az összes témakör.
A navigációt a lap tetején elhelyezett, minden oldalon elérhető fejléc biztosítja. 
A fejléc bal oldalán található a házikó ikon, amely a főoldalra vezet vissza.
 Mellette kettő lenyíló menü érhető el: 
 \begin{itemize}

    \item Számításelmélet I
    \item Számításelmélet II.

 \end{itemize}


A Számításelmélet I menü az első féléves anyagot, míg a Számításelmélet II menü a második félév témaköreit tartalmazza.
 A témakörök menün belüli sorrendje az előadások haladását követik,
 így a hallgatók a tananyag felépítésének megfelelően ezt szépen követve, járhatják be az egyes fejezeteket.

 A lenyíló listában pici pipa jelöli azokat az oldalakat,
  amelyeket a felhasználó már megnyitott.
 A menü alján egy számláló mutatja a haladást, „Megnyitva: 3 / 16”, tájékoztatást adva arról, hogy az adott menüből, (számításelmélet I. vagy II. csoportból ) hány oldala került már megtekintésre/feldolgozásra.
\begin{figure}[H] \centering \fbox{ \includegraphics[width=0.5\textwidth]{images/userdoksi_haladas.png} } \caption{A témakörválasztó menü haladáskövető funkciója} \label{fig:menu_haladas} \end{figure}


\newpage
A két menü által lefedett témakörök összefoglalása az alábbi táblázatban látható:


\begin{table}[H]
\centering
\begin{tabular}{p{0.46\textwidth} p{0.46\textwidth}}
\toprule
\textbf{Számításelmélet I} & \textbf{Számításelmélet II} \\
\midrule
Alapok 1: ábécé, szavak, műveletek        & Nulladrendű és elsőrendű logika \\
Alapok 2: nyelvek és műveleteik           & Turing-gépek \\
Generatív grammatikák                     & Számosság és eldönthetetlenség \\
Chomsky-hierarchia                        & Bonyolultsági osztályok \\
Lineáris grammatikák és a 3-as normálforma & NP-teljesség \\
Reguláris kifejezések                     & \\
Determinisztikus véges automata           & \\
Nemdeterminisztikus véges automata        & \\
Minimális automata és Myhill--Nerode tétel & \\
Reguláris kifejezések és automaták        & \\
Kis Bar-Hillel lemma                      & \\
Redukált grammatika és Chomsky-normálforma & \\
CYK algoritmus                            & \\
Veremautomata                             & \\
Nagy Bar-Hillel lemma                     & \\
Környezetfüggő és 0-típusú nyelvek        & \\
\bottomrule
\end{tabular}
\caption{Az alkalmazásban elérhető témakörök}
\label{tab:temakorok}
\end{table}

A fejlécben lévő ház ikonra kattintva a user bármikor, bármely oldalról, visszatérhet a főoldalra, így a navigáció a rendszer minden pontján egyszerűen elérhető.

\begin{figure}[H]
    \centering
    \includegraphics[width=0.7\textwidth]{images/userdoksi_fooldal1.png}
    \caption{Főoldal és megnyitott témakörválasztó menü (A Számításelmélet I)}
    \label{fig:főoldal menü1}
\end{figure}
% Mit mutass be:
%   - ragadós fejléc: FOTOC logó/ikon (home link), két lenyíló menü (Számításelmélet I / II), sötét–világos mód kapcsoló
%   - a lenyíló lista: témakörök neve, pipajel (✓) a már megnyitottaknál, "Megnyitva: X / Y" számláló alul
%   - a főoldal hero szövege (rövid leírás, mi a FOTOC, hogyan induljunk el)
% Kérdések: Milyen témakörök érhetők el? Hogyan lehet navigálni közöttük? Mit lát a felhasználó első megnyitáskor?
% Képernyőkép: főoldal + nyitott lenyíló menü
\mbox{}

\subsection{Egy témakör oldala}
% Mit mutass be:
%   - morzsaútvonal fejléc (pl. "Számításelmélet I > Reguláris nyelvek > DFA")
%   - nagy cím, alatta elméleti összefoglaló (TheoryRenderer rendereli: szöveg, képletek, definíciók)
%   - "Próbáld ki!" szekció – az interaktív modul (nem minden témakörnek van)
%   - "Teszteld magad!" szekció – kvízkérdések (nem minden témakörnek van)
% Kérdések: Milyen típusú tartalmak jelennek meg az elméleti részben? Mi történik ha nincs modul / nincs kvíz?
% Képernyőkép: egy teljes témakör oldal (pl. DFA), jelölve a három részt

A témakörök oldalai nagyjából mind egységes felépítést követnek.
Az oldal tetején egy úgynevezett morzsaútvonal (breadcrumb) segíti a tájékozódást,
 amely megmutatja, hogy a felhasználó a tananyag mely részén jár.
 
A fejlécben található navigációs elemek a témakörök oldalain is elérhetők.
 Segítségükkel a felhasználó bármikor visszatérhet a főoldalra, vagy átválthat egy másik témakörre.
Az oldal fő részét az elméleti összefoglaló alkotja. 
Ez tartalmazza az adott témakör megértéséhez szükséges definíciókat, tételeket, példákat, valamint matematikai képleteket és jelöléseket.

Az elméleti tartalom célja, hogy röviden összefoglalja az előadások legfontosabb fogalmait, és alapot biztosítson az interaktív feladatok használatához.

Egyes témakörök esetén az elméleti rész alatt egy „Próbáld ki!” szekció is megjelenik.

Ez egy interaktív szimulációt vagy vizualizációt tartalmaz, amely lehetővé teszi a tananyag valamilyen gyakorlati kipróbálását.
Egyes szimulációk saját példák megadását is lehetővé teszik, így a felhasználó a tananyag működését tetszőleges bemeneteken is kipróbálhatja, módosíthatja a bemenetet, és lépésről lépésre követheti az algoritmusok/automaták működését.

Azoknál a témaköröknél, ahol ellenőrző kérdések állnak rendelkezésre, az oldal alján egy „Teszteld magad!” szekció található. A kvízkérdések segítségével a hallgató azonnali visszajelzést kaphat arról, hogy mennyire sajátította el az adott témakör anyagát.
A kvízek elsősorban a szimulációval nehezebben szemléltethető, elméletibb témaköröknél jelennek meg
. Minden témakör legalább egy kiegészítő elemet tartalmaz, amely vagy interaktív szimuláció, vagy kvíz formájában jelenik meg.


A \ref{tab:temakorok}. táblázat összefoglalja az egyes témakörökhöz tartozó interaktív modulokat és kvízeket.
\begin{table}[H]
\centering
\small
\begin{tabular}{p{0.52\textwidth} p{0.26\textwidth} c}
\toprule
\textbf{Témakör} & \textbf{Interaktív modul} & \textbf{Kvíz} \\
\midrule
\multicolumn{3}{l}{\textit{Számításelmélet I}} \\
\midrule
Alapok 1: ábécé, szavak, műveletek        & Szójáték              & ---         \\
Alapok 2: nyelvek és műveleteik           & ---                    & \checkmark  \\
Generatív grammatikák                     & Grammatika levezető    & ---         \\
Chomsky-hierarchia                        & Chomsky-osztályozó    & ---         \\
Lineáris grammatikák és a 3-as normálforma & 3-as normálforma      & ---         \\
Reguláris kifejezések                     & Regex ellenőrző       & ---         \\
Determinisztikus véges automata           & DVA szimulátor        & ---         \\
Nemdeterminisztikus véges automata        & NVDA labor            & ---         \\
Minimális automata és Myhill--Nerode tétel & Minimalizáló          & ---         \\
Reguláris kifejezések és automaták        & ---                   & \checkmark  \\
Kis Bar-Hillel lemma                      & ---                   & \checkmark  \\
Redukált grammatika és Chomsky-normálforma & CF levezetés          & \checkmark  \\
CYK algoritmus                            & CYK vizualizátor      & ---         \\
Veremautomata                             & Veremautomata labor   & ---         \\
Nagy Bar-Hillel lemma                     & ---                   & \checkmark  \\
Környezetfüggő és 0-típusú nyelvek        & ---                   & \checkmark  \\
\midrule
\multicolumn{3}{l}{\textit{Számításelmélet II}} \\
\midrule
Nulladrendű és elsőrendű logika           & Logika labor          & ---         \\
Turing-gépek                              & Turing-gép szimulátor & ---         \\
Számosság és eldönthetetlenség            & ---                   & \checkmark  \\
Bonyolultsági osztályok                   & ---                   & \checkmark  \\
NP-teljesség                              & ---                   & \checkmark  \\
\bottomrule
\end{tabular}
\caption{Az alkalmazásban elérhető témakörök interaktív moduljaikkal és kvízeikkel}
\label{tab:temakorok}
\end{table}

\mbox{}

\subsection{Haladáskövetés és megjelenés} 
Az alkalmazás automatikusan nyomon követi, 
hogy a felhasználó melyik témaköröket nyitotta már meg. 
Egy témakör első megnyitását követően a navigációs menüben egy pipa jelenik meg az adott oldal mellett a lenyíló mneüben.
 A menü alján található számláló megmutatja, hogy az adott témakörcsoportból hány oldal került már megnyitásra.
  Ez segítséget nyújt a tananyag feldolgozásának követésében. 
  A fejléc jobb oldalán található nap/hold ikon, ennek segítségével a felhasználó bármikor válthat a világos és a sötét megjelenési mód között.
 A választott beállításokat és haladást az alkalmazás megjegyzi, ezért a következő megnyitáskor is megmaradnak. 
  A haladási adatok és a megjelenési mód beállításai az adott böngészőben kerülnek tárolásra. 
  Emiatt másik böngésző vagy eszköz használata esetén vagy incognito módban ezek az adatok
  nem kerülnek automatikusan átvitelre. 
  A tárolt adatok a böngésző helyi adatainak törlésével visszaállíthatók alaphelyzetbe.
\mbox{}

\section{Interakív modulok használata}
Az alkalmazás több témakörhöz interaktív modult biztosít.
 Ezek a modulok, az adott témakörhöz kapcsolódó automaták, algoritmusok és egyéb számításelméleti 
 modellek kipróbálását biztosítja.
A következő alfejezetekben az egyes interaktív modulok használata olvasható,
 beleértve a megadható bemeneteket, a vezérlési lehetőségeket és az eredmények értelmezését.
 
\newpage 
\subsection{Szóműveletek}
% Mit mutass be: (LanguageExplorer + WordPlayground komponensek)
%   - LanguageExplorer: automata vagy grammatika által generált/elfogadott szavak böngészése,
%     szűrés hossz szerint, példaszavak listázása
%   - WordPlayground: szabad szó-tesztelés – adott automatán vagy grammatikán próbálhat szavakat
% Kérdések: Véges vagy végtelen szólistát mutat? Hogyan lehet szűrni? Mi a különbség a két almodul között?
Az első és talán a legegyszerűbb modul a legelső, vagyis az "Alapok 1: ábécé, szavak, műveletek" témát egészíti ki. 
A modul az alábbi szóműveletek eredményét jeleníti meg egy megadott szó alapján: \begin{itemize} \item hossz ($|u|$), \item tükörkép ($u^R$), \item konkatenáció egy másik szóval ($uv$), \item hatványozás ($u^n$), \item prefixek és valódi prefixek, \item szuffixek és valódi szuffixek. \end{itemize}
Mindkét szó és a kitevő input szövegmezőkben adható meg az eredmény a \emph{Számol} gombra frissül.
 Az mintán az első szó \texttt{abc},
 a második szó \texttt{ba}, a kitevő pedig \texttt{2}. 
 A második szó csak a konkatenáció művelet eredményét befolyásolja. 
 \begin{figure}[H] \centering 
  \includegraphics[width=0.5\textwidth]{images/userdoksi_szomuveletek1.png} 
  \caption{Példa bemenet a szóműveletek modulban} 
  \label{fig:szomuveletek_input} 
\end{figure}

\begin{figure}[H]
    \centering
    \includegraphics[width=0.5\textwidth]{images/userdoksi_szomuveletek2.png}
\caption{A megadott bemenethez tartozó eredmények}
    \label{fig:főoldal menü1}
\end{figure}

\newpage
\subsection{Grammatikai levezetés}
% Mit mutass be: (GrammarModule = GrammarInput + GrammarDerivation)
%   - grammatika megadása: nemterminálisok, terminálisok, szabályok (pl. S → aSb | ε)
%   - levezetés lépésenkénti megjelenítése egy adott szóra
%   - három előre beállított példa: reguláris, CF és CS grammatika
% Kérdések: Milyen formátumban kell megadni a szabályokat? Mi történik ha a szó nem vezethető le?

A grammatikai levezetés modul a generatív grammatikák és a Redukált grammatika és Chomsky-normálforma témakörök gyakorlását segíti. A felhasználó saját nyelvtant adhat meg,
szerkesztése során megadhatók a szabályok, a terminális és nemterminális szimbólumok halmaza, valamint a kezdőszimbólum. 
A szabályokat \texttt{A -> x} formában kell megadni, a \texttt{->} baloldalán a kiinduló szímbolum, jobb oldalán pedig az a szimbólumsorozat szerepel, amelyre az adott nemterminális átírható. 
A levezetés során a felhasználó kiválaszt egy alkalmazható szabályt, 
majd megadja, hogy az aktuális szó melyik előfordulására szeretné alkalmazni.
 Az aktuális szó minden lépés után frissül. A 
 jobb oldali panelen a teljes levezetés gráfként követhető nyomon,
  ahol a csomópontok a közbülső szavakat jelölik. 
Ha az aktuális szó már csak terminális szimbólumokat tartalmaz, 
a modul jelzi, hogy a levezetés befejeződött.
A \emph{Vissza} gombbal az utolsó lépés visszavonható, az \emph{Újrakezdés} gombbal pedig a levezetés a kezdőszimbólumtól indítható újra.
 \begin{figure}[H] \centering 
  \includegraphics[width=1.0\textwidth]{images/userdoksi_grammatika.png} 
  \caption{Grammatika module} 
  \label{fig:szomuveletek_input} 
\end{figure}

\begin{figure}[H]
    \centering
    \includegraphics[width=0.5\textwidth]{images/userdoksi_grammatika2.png}
\caption{A levezetés menete}
    \label{fig:főoldal menü1}
\end{figure}



\subsection{Chomsky-osztályozó}
% Mit mutass be: (ChomskyClassifier komponens)
%   - grammatika bevitele (szabályok)
%   - automatikus osztályba sorolás (0–3-as típus) az összes Chomsky-feltétel ellenőrzésével
%   - visszajelzés: melyik feltétel teljesül / nem teljesül
% Kérdések: Milyen szabályformátumot fogad el? Megjeleníti-e a nem-teljesített feltételt?

A Chomsky-osztályozó modul egy megadott grammatikáról állapítja meg, hogy melyik Chomsky-típusba
 sorolható.

A grammatika megadásának formátuma azonos a grammatikai levezetés modulnál 
leírtakkal.
A beírt grammatikát a modul azonnal elemzi: 
minden egyes szabályhoz egy táblázatban mutatja meg, h
ogy az a 0-s, 1-es, 2-es és 3-as típus feltételeit teljesíti-e ($\checkmark$ jel), vagy sem (--).

Az 1-es típus (környezetérzékeny) esetén a modul a szabályokra rámutatva megjeleníti 
a talált $u_1 A u_2 \to u_1 v u_2$ felbontást is (az egér ráhúzásával látható).
Különleges esetet jelent a korlátozott $\varepsilon$-szabály (KES): 
ha van $S \to \varepsilon$ szabály, 
a modul külön ellenőrzi, hogy $S$ nem szerepel-e semmilyen szabály jobb oldalán,
 ha igen, a KES sérül, és a grammatika nem sorolható az 1-es típusba.
Ezt a vizsgálatot egy külön sávban, szöveges magyarázattal mutatja be.

Végeredményként a modul kiírja a grammatika típusát (pl. $\mathcal{G}_2$).

\begin{figure}[H]
    \centering
    \includegraphics[width=1.0\textwidth]{images/chomsky.png}
\caption{Chomsky module a példabemenetre}
    \label{fig:főoldal menü1}
\end{figure}



\subsection{3-as normálformára hozás} 

A modul egy környezetfüggetlen grammatikát alakít 3-as normálformára,
 miközben az átalakítás minden lépése nyomon követhető.
 A grammatika megadásának menete az előzőekben leírtakkal azonos.
 
 Az átalakítás két fő szakaszból áll: \begin{enumerate} \item \textbf{Hosszredukció:} a hosszú 
  jobb oldalak felbontása segédnemterminálisok segítségével. 
  \item \textbf{Láncmentesítés:} az $A \to B$ alakú egységszabályok eltávolítása. 
\end{enumerate} 

Minden lépéshez rövid magyarázat tartozik, az új szabályok kiemelve, a törölt szabályok áthúzva jelennek meg.
 A láncmentesítés során a program megjeleníti a kiszámított lánchalmazokat is.
A lejátszást a vezérlőgombokkal (elejére, vissza, lejátszás/megállítás, előre, végére)
 és a csúszkával lehet irányítani, a lejátszás sebessége is állítható (lassú/közepes/gyors).
\begin{figure}[H]
    \centering
    \includegraphics[width=0.7\textwidth]{images/normal3kezdo.png}
\caption{3-as normálformára hozó module}
    \label{fig:főoldal menü1}
\end{figure}
\begin{figure}[H]
    \centering
    \includegraphics[width=0.7\textwidth]{images/normal3form2.png}
\caption{3-as normálformára hozó module működés közben}
    \label{fig:főoldal menü1}
\end{figure}

\subsection{Reguláris kifejezés ellenőrző}

A Reguláris kifejezések kiegészítő modulja, előre elkészített feladatokon keresztül segíti a reguláris 
kifejezések gyakorlását.
 A felhasználónak a megadott nyelvekhez kell reguláris kifejezéseket készítenie.
A feladatok között a felület tetején található sorszámozott gombokkal lehet váltani.

A modulban az elméleti reguláris kifejezés szintaxisa érvényes.
 Az input szintaxis beírásához, a billentyűzet mellett gombok egészítik a kifejezések
bevitelét.
Az ellenőrzés eredménye háromféle lehet:
\begin{itemize} 
  \item \textbf{Helyes}, ha a két kifejezés ugyanazt a nyelvet írja le. 
  \item \textbf{Nagyon valószínű, hogy helyes}, ha a feladat szöveges nyelvleírás alapján csak korlátozott hosszúságú szavakig ellenőrizhető, és ezen a tartományon belül nem található eltérés. 
  \item \textbf{Ellenpélda}, ha a megadott kifejezés nem a megfelelő nyelvet írja le. Ilyenkor a rendszer kiadja a legrövidebb olyan szót, amelynél már eltérés mutatkozik. \end{itemize}
Hibás szintaxis esetén a rendszer hibaüzenetet jelenít meg. A feladatokhoz emellett tipp és egy lehetséges megoldás is megtekinthető.
\begin{figure}[H]
    \centering
    \includegraphics[width=0.7\textwidth]{images/regularis1.png}
\caption{Reguláris kifejezés ellenőrző modul: egy bemenet ellenőrzése}
    \label{fig:főoldal menü1}
\end{figure}


\subsection{Determinisztikus véges automata szimulátor}
% Mit mutass be: (DfaSimulator komponens)
%   - előre betöltött DFA (pl. paritásautomata {a,b} felett)
%   - szó megadása, lépésenkénti lejátszás (PlayerControls: lépés előre/hátra, lejátszás)
%   - az aktuális állapot kiemelése a gráfon (AutomatonGraph), átmenetek animálása
%   - elfogad/elutasít eredmény megjelenítése
% Kérdések: Hogyan adható meg a bemeneti szó? Mit jelent a zöld/piros kiemelés? Hány lépésig megy vissza?

A determinisztikus véges automata témakörhöz tartozó modul egy (DVA) szimulátora,
 amely az automata futását lépésenkénti animációval jeleníti meg egy user által állítható bemenetre.
A bemeneti szót egy szövegmezőbe kell beírni az elfogadott ábécé ($\Sigma$) a mező mellett olvasható.
Üres szó megadásához elegendő üresen hagyni a mezőt (az automata az $\varepsilon$ szóra is futtatható).
A szimulátorban a bemeneti szó betűi külön gombként jelennek meg, így a felhasználó tetszőleges lépésre 
ugorhat. Az állapotgráfon kiemelés jelzi a már bejárt útvonalat, az aktuális átmenetet és az aktuális 
állapotot. A lejátszásvezérlő gombokkal vagy a billentyűzet segítségével a 
szimuláció lépésenként vagy automatikusan is végigkövethető. Korábbi lépésekhez bármikor vissza lehet 
térni. A leolvasás végén a modul egyértelműen jelzi, hogy a szó elfogadásra vagy elutasításra került. 
Ábécén kívüli karakter esetén a szimuláció nem indítható el.
\begin{figure}[H]
    \centering
    \includegraphics[width=0.7\textwidth]{images/nemdet.png}
\caption{(DVA) szimulátor menet közben}
    \label{fig:főoldal menü1}
\end{figure}

\subsection{Nemdeterminisztikus véges automata témakör moduleja}
Ez a modul a nemdeterminisztikus véges automaták determinizálását szemlélteti a részhalmazos konstrukció
 segítségével. Az oldal tetején előre definiált 2 példa közül lehet választani, amelyek különböző eseteket 
 mutatnak be.
 \begin{itemize} 
  \item \textbf{Determinizálás}: egy általános NDVA példáján mutatja be a részhalmazos konstrukció működését, 
  beleértve az $\varepsilon$-átmenetek kezelését is. 
  \item \textbf{3NF-ből NDVA}: egy reguláris grammatikából felépített NDVA-t szemléltet, 
  bemutatva a reguláris nyelvtanok és a véges automaták közötti megfeleltetést. 
\end{itemize}

A user interakció a navigációra szűkül ennél a példák közüli válsztás után.
A lejátszást a vezérlőgombokkal (elejére, vissza, lejátszás/megállítás, előre, végére) és a csúszkával lehet irányítani,a sebessége is állítható (lassú/közepes/gyors).
A folyamat végén összefoglaló jelenik meg az elkészült DVA állapotainak számáról 
  és az esetleges csapdaállapotról.

\begin{figure}[H]
    \centering
    \includegraphics[width=0.7\textwidth]{images/nemdeter.png}
\caption{determinizáló module}
    \label{fig:főoldal menü1}
\end{figure}



\subsection{Minimalizáló}

A DFA-minimalizáló modul egy determinisztikus véges automata minimalizálásának folyamatát mutatja be lépésről lépésre. A modul egy előre betöltött automatával indul, amelynek átalakítását a felhasználó animált formában nézheti végig.
A levezetés lépései között a vezérlőgombokkal lehet navigálni.
 Lehetséges egyesével előre, hátra haladni, közvetlenül az első vagy utolsó lépésre ugrani,
  illetve a folyamat automatikus lejátszását elindítani. Az animáció sebessége pedig három fokozatban állítható.

Az algoritmus két fő szakaszból áll:

\begin{enumerate} \item \textbf{Elérhetőségi elemzés:} a modul megmutatja, hogy a kezdőállapotból mely állapotok érhetők el,
   majd eltávolítja a nem elérhető állapotokat.

\item \textbf{Particionálás:} a Myhill--Nerode-tételen alapuló algoritmus iteratív lépésekben finomítja
 az állapotok ekvivalenciaosztályait.
\end{enumerate}
A folyamat végén a modul megjeleníti az elkészült minimális automatát, 
amely az ekvivalenciaosztályokból létrehozott új állapotokat tartalmazza.
\begin{figure}[H]
    \centering
    \includegraphics[width=0.7\textwidth]{images/mini.png}
\caption{DFA-minimalizáló modul működés közben}
    \label{fig:főoldal menü1}
\end{figure}
\begin{figure}[H]
    \centering
    \includegraphics[width=0.7\textwidth]{images/minimegoldas.png}
\caption{DFA-minimalizáló modul eredménye}
    \label{fig:főoldal menü1}
\end{figure}

\subsection{CYK vizualizátor}

A CYK-vizualizátor a Cocke–Younger–Kasami algoritmus működését szemlélteti: megmutatja, hogy egy adott szó levezethető-e 
egy környezetfüggetlen grammatikából.

\textbf{Bemenet:}
\begin{itemize}
  \item \textbf{Grammatika:} az előző modulokhoz hasonló szerkesztőben adható meg. Fontos, hogy az algoritmus Chomsky-normálformát (CNF) vár, szóval amennyiben a grammatika nem CNF-alakú, a modul jelzi a sértő szabályt.
  \item \textbf{Vizsgált szó:} a jobb oldalon a grammatika input alatt, egy külön szövegmezőbe írható, üresen hagyva az $\varepsilon$ szóra fut.
\end{itemize}

Az algoritmus folyamatát és végül az eredménye egy háromszög alakú táblázatban jelenik meg:
 minden cella tartalmazza azon nemterminálisok halmazát, amelyek az adott részszóból levezethetők. A
  modul kiírja, hogy a teljes szó levezethető-e (vagyis a kezdőszimbólum szerepel-e a csúcscellában),
   és ha igen, megad egy lehetséges levezetést is.

Működés közben kiirásra kerül, hogy az adott részszóra melyik kettévágásokat próbálta ki az algoritmus,
egy vágásra víve az egeret látszik a hozzá tartozó két cella.
\begin{figure}[H]
    \centering
    \includegraphics[width=0.7\textwidth]{images/cykmuk.png}
\caption{CYK-vizualizátor module}
    \label{fig:főoldal menü1}
\end{figure}
\begin{figure}[H]
    \centering
    \includegraphics[width=0.7\textwidth]{images/cyk.png}
\caption{CYK-vizualizátor levezetés közben}
    \label{fig:főoldal menü1}
\end{figure}
\subsection{Veremautomata labor}

A modul egy konfigurálható veremautomatának szimulációját mutatja be, egy bementeti szóra futtatva
lépésenkénti animációval.

\textbf{Bal oldali konfigurációs panel:}
\begin{itemize}
  \item \textbf{Szabályok:} soronként egy átmenet adható meg a következő formátumban: \texttt{z q a -> u p}, ahol \texttt{z} a verem tetején lévő szimbólum, \texttt{q} az aktuális állapot, \texttt{a} a beolvasott betű (elhagyható $\varepsilon$-átmenetnél: \texttt{z q -> u p}), \texttt{u} a veremre írt szimbólumsorozat (balról jobbra, üres: \texttt{u} elhagyható), \texttt{p} a következő állapot.
  \item \textbf{Verem alja ($z_0$):} a verem aljaszimbóluma, alapértelmezetten \texttt{\#}.
  \item \textbf{Kezdőállapot ($q_0$)}: az automata induló állapota.
  \item \textbf{Végállapotok ($F$):} vesszővel elválasztva adhatók meg (végállapotos elfogadáshoz).
  \item \textbf{Elfogadás módja:} választható végállapotos ($L(A)$) vagy üres vermes ($N(A)$) elfogadás.
\end{itemize}

A panel alatt megjelenik az automatikusan leolvasott állapot, a veremszimbólum és ábécének a készletei,
 valamint egy jelzés arról, hogy az automata determinisztikus-e vagy nem.

A bemeneti szót a \emph{Bemeneti szó} mezőbe kell beírni,
 majd a \textbf{Futtat} gombra kattintani.
 
 A jobb oldali panelen az állapotgráf, az aktuális veremtartalom és az egyes lépések 
 részletei jelennek meg, és a lejátszásvezérlővel lépésenként 
 vagy automatikusan követhető a futás. 
 A leolvasás végén a modul jelzi, hogy a szó elfogadásra vagy elutasításra került.

\begin{figure}[H]
    \centering
    \includegraphics[width=0.7\textwidth]{images/vereminput.png}
\caption{Veremautomata bemeneti modulja, példabemenettel}
    \label{fig:főoldal menü1}
\end{figure}
\begin{figure}[H]
    \centering
    \includegraphics[width=0.7\textwidth]{images/cyk.png}
\caption{Veremaoutmata futtatás közben}
    \label{fig:főoldal menü1}
\end{figure} 

\subsection{Turing-gép szimulátor}

A Turing-gép szimulátor modul lehetővé teszi tetszőleges Turing-gépek futásának vizsgálatát. 
A felhasználó megadhatja a gép szabályait és a bemeneti szót.
A futást animálva tudja végigkövetni.

\textbf{Bal oldali konfigurációs panel:}
\begin{itemize}
  \item \textbf{Átmenetek:} soronként egy szabály adható meg: \texttt{állapot szimbólum -> állapot szimbólum irány}, ahol az irány \texttt{L} (bal), \texttt{R} (jobb) vagy \texttt{S} (helyben). A \texttt{\#} jellel kezdődő sorok megjegyzések.
  \item \textbf{Kezdőállapot ($q_0$):} az indulási állapot.
  \item \textbf{Elfogadó állapot ($q_i$):} ebbe az állapotba kerülve a gép elfogadja a szót.
  \item \textbf{Elutasító állapot ($q_n$):} ebbe az állapotba kerülve a gép elutasítja a szót.
  \item \textbf{Blank szimbólum ($\square$):} az üres cellák szimbóluma, alapértelmezetten \texttt{□}.
\end{itemize}

A panel alatt megjelenik az állapotok halmaza és az átmenetek száma. 
A bemeneti szót a \emph{Bemeneti szó} mezőbe kell beírni, 
majd a \textbf{Futtat} gombra kattintani. 
A jobb oldali panelen a szalag aktuális tartalma 
és a fej helyzete jelenik meg, az egyes lépések részletezésével. 
A lejátszásvezérlővel lépésenként vagy automatikusan követhető a futás. 
A gép megállása után egyértelműen jelzi, hogy elfogadta vagy elutasította-e a szót.

\subsection{Logika modul}

A logika modul egy inputként beadott logikai formula igazságtáblázatát építi fel lépésről lépésre, és lehetőséget ad a DNF, illetve KNF olvasat megjelenítésére.

\textbf{Bemenet:} a formula egy szövegmezőbe írható a következő operátorokkal:
\begin{itemize}
  \item tagadás: \texttt{\textasciitilde} vagy \texttt{!} (jelölés: $\lnot$)
  \item konjunkció: \texttt{\&} (jelölés: $\land$)
  \item diszjunkció: \texttt{|} (jelölés: $\lor$)
  \item implikáció: \texttt{->} (jelölés: $\to$)
  \item zárójelek csoportosításhoz
\end{itemize}
Példa: \texttt{y -> x \& \textasciitilde{}y | \textasciitilde{}z}.

A bal oldali panelen a formula szerkezeti fája jelenik meg,
 kiemelve az éppen kiértékelt részformulát. Az animáció a teljes fa megjelenítésétől indul, 
 majd oszloponként tölti fel az igazságtáblázatot.
  A lejátszásvezérlővel lépésenként vagy automatikusan követhető a kiértékelés. 
  Az elkészült táblázat alapján a modul kigombra megjeleníti a formula \textbf{DNF} és \textbf{KNF} alakját is.

\section{Kvízek}
Néhány témakörhöz  önellenőrző kvíz tartozik, amely segít a tananyag elsajátításának
 ellenőrzésében. 
 A kérdések két nehézségi szintre vannak bontva: 
 \textit{Könnyű} és \textit{Nehezebb} feladatokra. 
 A felhasználó feleletválasztós kérdések segítségével tesztelheti tudását, ahol 
 minden kérdéshez négy lehetséges válasz (A, B, C) tartozik.
A válasz kiválasztása után a rendszer azonnal visszajelzést ad. 
Helyes válasz esetén a kiválasztott lehetőség zöld színnel jelenik meg, 
hibás válasz esetén pedig piros színnel. Minden válasz után rövid magyarázat 
olvasható, amely segít megérteni a helyes vgay helytelen megoldás mögött álló elméleti hátteret.
A kvíz végén összesítő eredmény jelenik meg az elért pontszámmal. 
Az \textit{Újrakezdés} gomb segítségével a teszt bármikor újra kitölthető. 
Az aktuális állapot és az eredmények a böngésző helyi tárhelyén tárolódnak, 
így az oldal későbbi megnyitásakor a korábban elért eredmények visszatölthetők.
\begin{figure}[H]
    \centering
    \includegraphics[width=0.7\textwidth]{images/kviz.png}
\caption{Kvízek kiértékelése}
    \label{fig:főoldal menü1}
\end{figure}
\begin{figure}[H]
    \centering
    \includegraphics[width=0.7\textwidth]{images/kviz2.png}
\caption{Kvízek az összesítő eredménye és az Újrakezdés gomb }
    \label{fig:főoldal menü1}
\end{figure} 




\section{Jelölések}

Az alkalmazásban vannak olyan beviteli szintaxisok, amik eltérnek a papíros jelöléstől. Az alábbi táblázat összefoglalja a főbb konvenciókat.

\begin{table}[H]
\centering
\begin{tabular}{p{0.30\textwidth} p{0.30\textwidth} p{0.30\textwidth}}
\toprule
\textbf{Fogalom} & \textbf{Papíros jelölés} & \textbf{Beviteli formátum} \\
\midrule
Üres szó & $\varepsilon$ & üres mező (semmit sem kell beírni) \\
Nemterminálisok & $\{S, A, B\}$ & \texttt{S,A,B} (vesszővel elválasztva) \\
Terminálisok & $\{a, b\}$ & \texttt{a,b} (vesszővel elválasztva) \\
Grammatikaszabály & $A \to BC$ & \texttt{A->BC} (soronként egy szabály) \\
$\varepsilon$-szabály & $S \to \varepsilon$ & \texttt{S->} (jobb oldal üres marad) \\
Turing-átmenet & $\delta(q_0, a) = (q_1, b, R)$ & \texttt{q0 a -> q1 b R} (soronként) \\
Logikai tagadás & $\lnot p$ & \texttt{\textasciitilde{}p} vagy \texttt{!p} \\
Logikai konjunkció & $p \land q$ & \texttt{p \& q} \\
Logikai diszjunkció & $p \lor q$ & \texttt{p | q} \\
Logikai implikáció & $p \to q$ & \texttt{p -> q} \\
\bottomrule
\end{tabular}
\caption{Beviteli konvenciók az alkalmazásban}
\label{tab:jeloelesek}
\end{table}

Az állapotok neve tetszőleges szöveges azonosító lehet (pl.\ \texttt{q0}, \texttt{q\_accept}, \texttt{START}). A Turing-gép átmenetibeállító mezőjében a \texttt{\#} karakterrel kezdődő sorok megjegyzésnek számítanak és figyelmen kívül maradnak. Az irány értéke \texttt{L} (bal), \texttt{R} (jobb) vagy \texttt{S} (helyben) lehet.

Az elméleti összefoglalókban és egyes feladatokban megjelenő matematikai képleteket a rendszer KaTeX segítségével rendereli,
 ezért azok a tankönyvekben megszokott matematikai formában jelennek meg.

\section{Hibaüzenetek és hibaelhárítás}

\subsection{Tipikus beviteli hibák}

Az alkalmazás a legtöbb beviteli hibát azonnal, az érintett mező alatt jelzi.

\begin{itemize}
  \item \textbf{Érvénytelen grammatikaszabály:} ha egy szabály bal vagy jobb oldala olyan szimbólumot tartalmaz, amely nincs felsorolva a nemterminálisok vagy terminálisok között, az alkalmazás sárga figyelmeztető szöveget jelenít meg (pl.\ \textit{„\texttt{S->XY}: jobb oldal ismeretlen szimbólumot tartalmaz"}).
  \item \textbf{Grammatika nem CNF-alakú (CYK modul):} ha a beírt grammatika nem Chomsky-normálformában van, piros keretes hibaüzenet jelenik meg, megnevezve a sértő szabályt és az okot.
  \item \textbf{Helytelen Turing-átmenet:} ha egy sor nem tartalmazza a \texttt{->} elválasztót, vagy az irány érvénytelen, a modul sárga szöveggel jelzi az érintett sor számát és a várt formátumot (\textit{„Helytelen formátum -- várt: állapot szimbólum -> állapot szimbólum irány"}).
  \item \textbf{Üres kötelező mező:} ha a Turing-gépen a kezdő-, elfogadó vagy elutasító állapot nincs kitöltve, a „Futtat" gomb inaktív marad.
\end{itemize}

\subsection{Ismert korlátok}

\begin{itemize}
  \item A CYK modul legfeljebb \textbf{10 karakteres} bemeneti szót fogad el, mivel ennél hosszabb szavaknál a háromszög-táblázat áttekinthetetlenné válik.
  \item Nagyon sok átmenetet tartalmazó automatáknál a gráf megjelenítése lassabb lehet.
  \item Az alkalmazás mobileszközön korlátozottan használható: a kisebb képernyőn egyes interaktív elemek (pl.\ a gráf mozgatása, az átmenettábla) szűk területen jelennek meg.
\end{itemize}

\subsection{Ha az oldal nem töltődik be}

Online verzió esetén ellenőrizze az internetkapcsolatot, és töltse újra az oldalt. Helyi futtatásnál ellenőrizze, 
hogy a \texttt{pnpm --filter web dev} parancs hiba nélkül elindult-e, és a \texttt{http://localhost:5173} cím elérhető-e a böngészőből. 
A szükséges szoftverek verziójára vonatkozó előírások a \ref{ch:user}. fejezet elején vannak felsorolva.

\subsection{Hibajelentés}

Hibák és javaslatok jelezhetők a projekt GitHub-oldalán az Issues felületen:
\begin{center}
  \url{https://github.com/SZaboAnna21/thesis/issues}
\end{center}
